\documentclass[11pt]{article}
\usepackage[T1]{fontenc}
\usepackage{textcomp}
\usepackage[margin=0.75in]{geometry}
\usepackage{amsmath,amssymb,amsthm}
\usepackage{array,booktabs}
\usepackage{hyperref}
\setlength{\parindent}{0pt}
\newtheorem{theorem}{Theorem}
\theoremstyle{definition}
\newtheorem{definition}{Definition}
\title{Flow Proof Artifact: prop-02-perpendicular-bisects-chord}
\author{Generated by \texttt{flow doc proof}}
\date{Flow Proof Book}
\begin{document}
\maketitle
If a straight line through the center of a circle cuts a chord at right angles, it bisects the chord and the central angles.\\[0.5em]
\textbf{Source.} euclid — Elements, Book III, Proposition 2\\[0.5em]
\bigskip
\noindent{\large \textbf{Derived fact 1.} \textit{If a straight line through the center of a circle cuts a chord at right angles, it bisects the chord and the central angles} \hfill the Euclidean plane \cdot Euclid Book III \cdot Proposition 2: a perpendicular from the center bisects the chord}\par
\smallskip\noindent\textit{Coordinate: the Euclidean plane \cdot Euclid Book III \cdot Proposition 2: a perpendicular from the center bisects the chord}\par
\smallskip\noindent\textit{Source: euclid — Elements, Book III, Proposition 2}\par
\smallskip\noindent\textit{Built on: proposition 10: bisect a given segment, for Book I of the Elements in on the Euclidean plane}\par
\medskip
\noindent\textbf{Goal.} If a straight line through the center of a circle cuts a chord at right angles, it bisects the chord and the central angles\par
\noindent$\displaystyle \text{perpendicular from center bisects chord}$\par
\medskip
\noindent
\begin{tabular}{@{} >{\raggedright\arraybackslash}p{0.06\textwidth} >{\raggedright\arraybackslash}p{0.47\textwidth} >{\raggedleft\arraybackslash}p{0.41\textwidth} @{}}
\textbf{\#} & \textbf{Proof} & \textbf{Mathematics} \\
\midrule
\textbf{1.} & We prove that proposition 2: a perpendicular from the center bisects the chord for Euclid Book III the Euclidean plane. & \\[0.45em]
\textbf{2.} & Let O = center. & \\[0.45em]
\textbf{3.} & Let AB = chord. & \\[0.45em]
\textbf{4.} & Let M = foot\_of\_perpendicular. & \\[0.45em]
\textbf{5.} & From step 2, step 3, and step 4, we can deduce that segment AM equals segment MB. & $\displaystyle \text{segment AM} = \text{segment MB}$ \\[0.45em]
\textbf{6.} & We invoke the derived fact governing Book I of the Elements in the Euclidean plane: proposition 10: bisect a given segment, for Book I of the Elements in on the Euclidean plane. & \\[0.45em]
\textbf{7.} & From step 6, this implies perpendicular from center bisects chord. Hence proven. & $\displaystyle \text{perpendicular from center bisects chord}$ \\[0.45em]
\bottomrule
\end{tabular}
\label{thm:Geometry-Euclid Book III-proposition_2_a_perpendicular_from_the_center_bisects_the_chord}
\par\medskip\hrule
\end{document}
